% Ek C -- KITTI ek sonuçları (revizyon-v2: özgün Ek A'nın KITTI kısmı + ana metindeki KITTI kurulum ayrıntıları,
% best.pt artefakt analizi (eski 8.6(b)), H07 KITTI ayrıntıları ve gecikme hatası (eski 8.6(c))).
% main.tex ve cevrimici_ek.tex \input eder.
\setcounter{table}{0}\setcounter{figure}{0}
\section{KITTI ek sonuçları}\label{sec:ek-ikincil}

Bu ek, KITTI deneyinin (Bölüm~\ref{sec:genelleme}) kurulum ayrıntılarını, ikincil sonuçlarını, H05'in eğitim eğrisini, \texttt{best.pt} seçimi analizini, kilit sonrası sağlamlık analizini ve eğitimsiz çıkarım ayrıntılarını verir.

\subsection{Kurulum ayrıntıları}\label{sec:ek-kitti-kurulum}

\paragraph{Sorular, kurallar ve eğitim ayrıntıları.}
S1 model başına iki aileyle yanıtlanır: iyileşme ailesinde kabul ölçütünü sağlayan kart sayısı ve yönlü replikasyon ailesinde VOC'taki kaybın yinelenip yinelenmediği. S2, VOC'ta koşulan iki bileşen kartının katkı ve toplamsallık ölçütleriyle yanıtlanır. Kural belgesindeki sonraki değişiklikler S2'nin VOC kartlarına uygulanmaz. VOC'taki fiilî optimizatör davranışı şu ayarlarla sabitlendi: MuSGD, öğrenme hızı 0.01, momentum 0.937, bias ısınması 0. Copy-paste parametreleri değişmedi; mosaic başına yapıştırılan nesne H03'te 0.98, H08'de 2.00'dir. YOLO26n'de H02 ve H05 yamaları iki eğitim başlığına da uygulanır (başlık sayacıyla denetlendi). Üç tohumla Holm-5'in ilk adımında \%50 güçteki kritik etki $\approx2.8\,s_\Delta$'dır; kabul edilmeyen kart için ``kabul ölçütü sağlanmadı'' denir. H07'de NMS'siz YOLO26n'de NMS, Soft-NMS ve WBF satırları uygulanamaz.

\paragraph{Bölme ve yok-say kuralı.}
Hazır Ultralytics bölmesi dışında, 3B tespit çalışmalarında kullanılan 3712/3769 eğitim/doğrulama bölmesi de kullanılmadı: ayrı bir test bölümü sunmaz ve aynı sahnenin farklı kayıt adlarıyla gelen karelerini birleştiren bir denetim (NCC) içerdiğini doğrulayamadık. Sahne grupları eğitim, doğrulama ve test bölümlerine, her sınıfın nesnelerinin en az \%40, \%5 ve \%10'u ilgili bölmede olacak ve tek bir grup doğrulama ya da test hedef büyüklüğünün \%25'ini aşmayacak biçimde atandı. Bölmeler arası en yüksek NCC test--eğitim için 0.855, doğrulama--eğitim için 0.888'dir. Test bölümü 6\,212 nesne ve 2\,436 \texttt{DontCare} bölgesi içerir. Değerlendirmede \texttt{DontCare} bölgeleri her sınıf için COCO tarzı bir yok-say bölgesi (\texttt{iscrowd}) olarak eklenir. Tespitler önce gerçek nesnelerle eşleşir; yalnız eşleşmeyen tespit, bir \texttt{DontCare} bölgesiyle örtüşme oranı (IoA: kesişimin tespit alanına oranı) o anki IoU eşiğine eşit ya da büyükse yok sayılır. Resmî KITTI AP'si ise sınıf başına sabit eşik ve zorluk düzeyleri kullanır. \texttt{DontCare}'siz mAP ikincil olarak Tablo~\ref{tab:kitti-ikincil}'de verilir.

\paragraph{Eğitim, sağlık kapısı ve öğrenme denetimi.}
YOLO26n'nin kayıp kodu v8.4.165 \texttt{E2ELoss}'tur. KITTI'deki yineleme sayısı VOC'takinin (24\,300) $\approx$3.1 kat altındadır. Başlık sayacı her H02 ve H05 koşumunda okundu. Bir eğitim koşumu YOLO11s için 1.68--1.83, YOLO26n için 1.17--1.31 GPU saati sürdü. Her koşum otomatik bir kapıdan geçer: batch'in sessizce küçültülmemesi, yamanın gerçekten çağrılması ve bir öğrenme denetimi. Öğrenme denetimi geçerlilik ölçütüdür: \texttt{results.csv} tam 100 epoch içerir, kayıp ve metrik değerleri sonludur, öğrenme hızı planı aynı modelin tohum-0 taban modeliyle epoch epoch aynıdır, 81--90. epoch'ların ortalama eğitim (kutu $+$ sınıflandırma) kaybı 11--20. epoch'larınkinden düşüktür ve en iyi doğrulama mAP50-95 0.02'nin üstündedir. Taban modele göreli bir eşik konmadı, çünkü güçlü negatif kartları sansürlerdi. Geçmeyen koşum $\Delta$'ya girmez. En iyi epoch'un 20'den küçük olması geçersizlik değil, uyarıdır; kural tüm kartlara ve iki modele aynı biçimde uygulanır. Bu denetimin kural belgesine girişi ve test betiğine \texttt{last.pt} değerlendirmesinin eklenmesi kural değişiklikleridir (Bölüm~\ref{sec:sinirliliklar}). YOLO11s kuyruğu bitene kadar izleme yalnız olay türü, koşum kimliği ve çıkış koduyla yapıldı.

\paragraph{Çözünürlük.}
KITTI karelerinin çoğu 1242$\times$375'tir (testte 1\,396/1\,504; en-boy oranı 3.31) ve 640 px girişte uzun kenar $\approx$0.52 ölçeğine inerek kare $\approx$640$\times$193 olur. Testte pedestrian ve cyclist nesnelerinin COCO ``küçük'' sınıfında (alan $<32^2$ px) olanlarının oranı özgün çözünürlükte \%19.1, 640 px'te \%65.6'dır. VOC görüntüleri ise en çok 500 px'tir ve 640 px'te 1.28 kat büyütülür (VOC testte küçük nesne oranı \%4.7'den \%1.85'e iner). Eğitim çözünürlüğü tek değişken kuralı gereği VOC'la aynı (640 px) tutuldu. 832 px test girişinin KITTI'de kazandırıp VOC'ta kaybettirmesinin ve çok ölçekli eğitimin (H01) KITTI'de kabul alıp VOC'ta almamasının en basit hipotezi bu farktır: KITTI'de küçük nesne payı yüksek olduğundan daha yüksek etkin çözünürlük yarar sağlayabilir, VOC'ta nesneler 640 px'te zaten büyütülmüştür. Hipotez sınanmadı.

\paragraph{Taban model makullüğü.}
KITTI'de bu bölme ve değerlendirme kuralıyla yayımlanmış bir çıpa bulunmadığından VOC'taki makullük kapısı uygulanmadı. Yerine üç denetim yapıldı. Birincisi eğitim hattı tanısıdır: taban model \texttt{last.pt} ağırlıkları eğitim görüntülerinde yüksek, doğrulamada düşük mAP50-95 verir (sayılar: Ek~\ref{sec:ek-kontrol-noktasi}). İkincisi mutlak değerlerdir (\texttt{DontCare} kurallı, üç tohum ortalaması; puan): test mAP50-95'i YOLO11s'te 19.30 (\texttt{best.pt}) ve 19.66 (\texttt{last.pt}), YOLO26n'de 18.85 ve 17.35; \texttt{best.pt} doğrulama mAP50-95'i YOLO11s'te 19.26 $\pm$ 2.10, YOLO26n'de 18.51 $\pm$ 1.49; test car AP50-95'i (\texttt{best.pt}) YOLO11s'te 57.66, YOLO26n'de 58.10'dur. Bu değerler yalnız $\Delta$'ların ölçeğini verir; kıyas çıpası değildir, VOC değerleriyle ya da modeller arasında kıyaslanmaz. Üçüncüsü, yalnız büyüklük mertebesi için bir dış kaynak değeridir (Ek~\ref{sec:ek-dis}); çözünürlük, bölme ve değerlendirme farklı olduğundan bizim değerlerimizle karşılaştırılamaz.

\subsection{KITTI ikincil sonuçları ve koşum başına en iyi epoch}
Bu alt bölüm, KITTI kural belgesinde tanımlı ikincil raporları verir. Hepsi betimseldir; Holm düzeltmesi ve anlamlılık iddiası yoktur.

% uretici: evidence/scripts/build_ek_tables.py <- evidence/evidence.json (ek.test.*, ek.val.*, ek3.test.*, ek3.b2.*; elle duzenleme bir sonraki uretimde kaybolur)
\begin{table}[!htbp]
\centering\small\setlength{\tabcolsep}{6pt}
\caption{KITTI önceden belirlenmiş ikincil betimsel kıyaslar: kart $-$ taban model $\Delta$ ortalama $\pm$ s.d. [puan], $n=3$ tohum-eşli; Holm yok, renk yok.}\label{tab:kitti-ikincil}
\begin{tabularx}{\textwidth}{>{\raggedright\arraybackslash}X c c c c}
\toprule
\multicolumn{1}{@{}l}{} & \multicolumn{3}{c}{Test $\Delta$} & Val $\Delta$ \\
\cmidrule(lr){2-4}
Kart & DontCare'siz & 3 sınıf & 7 sınıf & ep11--90 eğri \\
\midrule
\multicolumn{5}{@{}l}{\textbf{YOLO26n}} \\
H01 multi-scale & +0.05 $\pm$ 2.19 & +1.22 $\pm$ 1.18 & +0.08 $\pm$ 2.52 & $-$0.31 $\pm$ 0.03 \\
H02 Inner-CIoU & $-$1.01 $\pm$ 1.98 & +1.04 $\pm$ 1.17 & $-$1.02 $\pm$ 2.29 & +0.51 $\pm$ 0.52 \\
H03 kutu copy-paste & +0.88 $\pm$ 0.68 & +2.81 $\pm$ 0.65 & +1.07 $\pm$ 0.87 & +0.62 $\pm$ 1.02 \\
H04 mosaic scale & $-$1.33 $\pm$ 2.19 & $-$0.97 $\pm$ 1.38 & $-$1.53 $\pm$ 2.46 & $-$0.54 $\pm$ 0.48 \\
H05 Varifocal & +0.47 $\pm$ 1.48 & +1.28 $\pm$ 1.16 & +0.78 $\pm$ 1.73 & +1.29 $\pm$ 0.20 \\
H06 kutu kazancı & $-$0.24 $\pm$ 0.72 & +0.56 $\pm$ 1.68 & $-$0.26 $\pm$ 0.86 & +0.33 $\pm$ 0.71 \\
H08 yüksek doz & +0.67 $\pm$ 0.68 & +2.68 $\pm$ 1.32 & +0.73 $\pm$ 0.89 & +1.01 $\pm$ 0.83 \\
\midrule
\multicolumn{5}{@{}l}{\textbf{YOLO11s}} \\
H01 multi-scale & +1.48 $\pm$ 0.39 & +0.48 $\pm$ 1.72 & +1.72 $\pm$ 0.29 & +0.59 $\pm$ 0.57 \\
H02 Inner-CIoU & +1.13 $\pm$ 0.48 & +0.06 $\pm$ 1.75 & +1.38 $\pm$ 0.66 & $-$0.27 $\pm$ 0.55 \\
H03 kutu copy-paste & +1.17 $\pm$ 1.00 & +0.22 $\pm$ 2.73 & +1.50 $\pm$ 1.17 & +1.01 $\pm$ 0.55 \\
H04 mosaic scale & +1.00 $\pm$ 0.83 & $-$0.66 $\pm$ 1.86 & +0.86 $\pm$ 1.30 & $-$0.36 $\pm$ 1.36 \\
H05 Varifocal & +2.35 $\pm$ 0.96 & +1.41 $\pm$ 2.06 & +2.73 $\pm$ 1.60 & +0.53 $\pm$ 0.37 \\
H06 kutu kazancı & +0.45 $\pm$ 0.65 & +0.13 $\pm$ 1.31 & +0.57 $\pm$ 0.79 & $-$0.38 $\pm$ 0.29 \\
H08 yüksek doz & +1.63 $\pm$ 2.89 & +1.30 $\pm$ 4.20 & +1.66 $\pm$ 4.04 & +1.00 $\pm$ 0.39 \\
\bottomrule
\addlinespace[3pt]
\multicolumn{5}{@{}p{\dimexpr\textwidth-2\tabcolsep\relax}@{}}{Test (\texttt{best.pt}): DontCare'siz mAP50-95; DontCare kurallı alt kümeler: car, pedestrian, cyclist (3 sınıf) ve person\_sitting hariç 7 sınıf. Val: eğitim eğrisinin (Ultralytics, DontCare'siz) 11.--90. epoch ortalaması. Test nesne sayısı: person\_sitting 56, tram 126, truck 119.} \\
\end{tabularx}
\end{table}


Tablo~\ref{tab:kitti-ikincil}'de \texttt{DontCare}'siz test $\Delta$'sı, 14 kartın hiçbirinde birincil (\texttt{DontCare} kurallı) $\Delta$'dan 0.24 puandan fazla ayrılmaz; birincil sonuç yok-say kuralına duyarlı değildir. YOLO11s H01'in kazancı standart üç sınıfta (car, pedestrian, cyclist) $+0.48\pm1.72$'dir (tohum $\Delta$'ları +1.57, $-$1.50, +1.38); person\_sitting hariç yedi sınıflık alt kümede $+1.72\pm0.29$'dur (üç tohumda pozitif). Fark, standart üç sınıfa van, truck, tram ve misc eklenince ortaya çıkar; bu sınıfların test nesne sayıları car'dan (4\,457) çok azdır (van 464, misc 184, tram 126, truck 119; person\_sitting 56). Bu sayılarla sınıf başına AP gürültülüdür ve kazancın tek tek sınıflara dağılımı bu veriyle ayrıştırılamaz. Ters yönde bir örnek: YOLO26n'de H03'ün üç sınıflık farkı $+2.81\pm0.65$ (üç tohumda pozitif), H08'inki $+2.68\pm1.32$'dir; bu iki kartın birincil $\Delta$'ları bunun çok altındadır (Tablo~\ref{tab:kitti-test}). Doğrulama eğri sütunu, eğitim sırasındaki doğrulama mAP50-95'inin 11.--90. epoch ortalamasının taban modelden farkıdır; H05 için YOLO26n'de $+1.29\pm0.20$, YOLO11s'te $+0.53\pm0.37$'dir (Ek~\ref{sec:ek-h05-kitti}).

% uretici: evidence/scripts/build_ek_tables.py <- evidence/evidence.json (ek.epoch.*, ek3.val.*; elle duzenleme bir sonraki uretimde kaybolur)
\begin{table}[!htbp]
\centering\small\setlength{\tabcolsep}{4pt}
\caption{KITTI: koşum başına en iyi epoch (\texttt{best.pt}) ve val best$-$last farkı [puan]; 48 koşum, tohum s0--s2. Eğik: epoch 1.}\label{tab:kitti-epoch}
\begin{tabularx}{\textwidth}{>{\raggedright\arraybackslash}X c c c c c c}
\toprule
\multicolumn{1}{@{}l}{} & \multicolumn{2}{c}{s0} & \multicolumn{2}{c}{s1} & \multicolumn{2}{c}{s2} \\
\cmidrule(lr){2-3}\cmidrule(lr){4-5}\cmidrule(lr){6-7}
Kol & Epoch & Fark & Epoch & Fark & Epoch & Fark \\
\midrule
\multicolumn{7}{@{}l}{\textbf{YOLO26n}} \\
B taban model & 16 & +4.3 & 11 & +5.1 & 28 & +3.7 \\
H01 multi-scale & 24 & +3.8 & 9 & +4.6 & 20 & +5.3 \\
H02 Inner-CIoU & 29 & +3.8 & 11 & +4.0 & 38 & +3.1 \\
H03 kutu copy-paste & 10 & +4.7 & 24 & +4.4 & 37 & +2.6 \\
H04 mosaic scale & 16 & +5.5 & 11 & +4.0 & 12 & +5.4 \\
H05 Varifocal & 46 & +2.9 & 39 & +4.2 & 33 & +2.8 \\
H06 kutu kazancı & 51 & +2.9 & 11 & +4.8 & 28 & +3.4 \\
H08 yüksek doz & 39 & +3.4 & 12 & +4.6 & 48 & +3.3 \\
\midrule
\multicolumn{7}{@{}l}{\textbf{YOLO11s}} \\
B taban model & 9 & +4.4 & 32 & +3.1 & \textit{1} & +5.7 \\
H01 multi-scale & 15 & +4.6 & 48 & +3.9 & \textit{1} & +5.5 \\
H02 Inner-CIoU & 29 & +4.6 & 38 & +4.6 & \textit{1} & +6.6 \\
H03 kutu copy-paste & \textit{1} & +4.8 & 12 & +2.7 & \textit{1} & +5.2 \\
H04 mosaic scale & 39 & +5.1 & \textit{1} & +7.4 & \textit{1} & +6.1 \\
H05 Varifocal & 47 & +2.7 & 54 & +1.9 & 17 & +3.2 \\
H06 kutu kazancı & 13 & +4.9 & \textit{1} & +3.6 & \textit{1} & +6.5 \\
H08 yüksek doz & \textit{1} & +2.9 & \textit{1} & +6.3 & 24 & +3.9 \\
\bottomrule
\addlinespace[3pt]
\multicolumn{7}{@{}p{\dimexpr\textwidth-2\tabcolsep\relax}@{}}{Val farkı: eğitim-içi en iyi $-$ son epoch mAP50-95 (Ultralytics, DontCare'siz). En iyi epoch 1 olan koşum: YOLO26n 0/24, YOLO11s 11/24.} \\
\end{tabularx}
\end{table}


Tablo~\ref{tab:kitti-epoch} her koşumun en iyi epoch'unu ve eğitim-içi en iyi eksi son epoch doğrulama farkını verir. YOLO11s'te 24 koşumun 11'inde en iyi epoch 1'dir (tohum 0'da 2, tohum 1'de 3, tohum 2'de 6); YOLO26n'de hiçbirinde değildir. Fark 48 koşumun tümünde pozitiftir (\texttt{best.pt} doğrulama üzerinde seçildiğinden beklenir) ve Ultralytics'in eğitim-içi (\texttt{DontCare}'siz) değerlendiricisiyle ölçülür: 24 koşum ortalaması YOLO26n'de $+4.03$, YOLO11s'te $+4.58$, yalnız taban model için $+4.37$ ve $+4.39$ puandır. Bu sayılar, Ek~\ref{sec:ek-kontrol-noktasi}'teki $+3.23$ ve $+3.84$'ten (bu çalışmanın COCO-tarzı yeniden değerlendirmesi) farklı bir ölçümdür.

% uretici: evidence/scripts/build_ek_tables.py <- evidence/evidence.json (ek4.test.*, ek4.val.*, ek.val.*; elle duzenleme bir sonraki uretimde kaybolur)
\begin{table}[!htbp]
\centering\small\setlength{\tabcolsep}{3pt}
\caption{KITTI ikincil ölçütler (\texttt{best.pt}): kart $-$ taban model $\Delta$ ortalama $\pm$ s.d. [puan], $n=3$ tohum-eşli; test DontCare kurallı. Betimsel, renk yok.}\label{tab:kitti-test-ek}
\begin{tabularx}{\textwidth}{>{\raggedright\arraybackslash}X c c c c c c}
\toprule
\multicolumn{1}{@{}l}{} & \multicolumn{4}{c}{Test $\Delta$} & \multicolumn{2}{c}{Val [puan]} \\
\cmidrule(lr){2-5}\cmidrule(lr){6-7}
Kart & mAP50 & AP$_S$ & AP$_M$ & AP$_L$ & $\Delta$ mAP50-95 & İç$-$yeniden \\
\midrule
\multicolumn{7}{@{}l}{\textbf{YOLO26n}} \\
B & 33.59$\pm$0.89 & 7.92$\pm$2.61 & 15.94$\pm$1.05 & 32.71$\pm$0.94 & -- & +0.35$\pm$1.20 \\
H01 & +0.23$\pm$1.82 & +0.05$\pm$1.21 & $-$0.37$\pm$1.48 & +2.32$\pm$3.38 & $-$0.10$\pm$2.25 & +0.21$\pm$1.01 \\
H02 & $-$1.31$\pm$2.54 & +0.29$\pm$2.68 & $-$0.51$\pm$2.18 & $-$1.86$\pm$0.40 & $-$0.29$\pm$0.80 & +0.51$\pm$0.52 \\
H03 & +1.94$\pm$0.43 & $-$0.10$\pm$1.47 & +1.05$\pm$0.18 & +1.81$\pm$3.11 & +0.63$\pm$1.06 & +0.19$\pm$0.46 \\
H04 & $-$1.90$\pm$2.14 & $-$1.80$\pm$1.65 & $-$0.83$\pm$2.08 & $-$1.38$\pm$0.97 & $-$0.22$\pm$1.89 & +0.78$\pm$0.59 \\
H05 & +0.40$\pm$2.14 & +0.13$\pm$1.82 & +1.02$\pm$1.75 & +0.57$\pm$2.08 & +0.68$\pm$1.14 & +0.44$\pm$0.53 \\
H06 & $-$0.38$\pm$1.10 & $-$0.30$\pm$1.92 & +0.18$\pm$1.31 & $-$0.90$\pm$1.91 & +0.09$\pm$0.65 & +0.20$\pm$0.32 \\
H08 & +1.18$\pm$0.69 & $-$1.24$\pm$3.08 & +0.72$\pm$1.54 & +0.96$\pm$1.53 & +0.46$\pm$1.34 & +0.29$\pm$0.52 \\
\midrule
\multicolumn{7}{@{}l}{\textbf{YOLO11s}} \\
B & 35.39$\pm$0.65 & 7.20$\pm$0.91 & 18.14$\pm$0.80 & 31.92$\pm$0.86 & -- & +1.12$\pm$0.68 \\
H01 & +1.63$\pm$1.08 & +2.58$\pm$1.61 & +0.98$\pm$1.77 & +2.74$\pm$2.76 & +1.73$\pm$1.78 & $-$0.01$\pm$0.33 \\
H02 & +1.87$\pm$2.46 & +1.28$\pm$0.84 & +0.95$\pm$0.84 & +2.10$\pm$1.24 & +1.15$\pm$1.05 & +0.28$\pm$0.29 \\
H03 & +1.84$\pm$2.56 & +2.50$\pm$2.35 & +0.60$\pm$1.22 & +1.28$\pm$1.94 & +1.86$\pm$2.44 & +0.65$\pm$0.54 \\
H04 & +2.14$\pm$0.26 & +2.81$\pm$3.12 & +0.81$\pm$1.34 & +2.50$\pm$1.03 & +1.82$\pm$0.43 & +0.74$\pm$1.00 \\
H05 & +1.66$\pm$0.76 & +4.37$\pm$2.90 & +1.60$\pm$1.01 & +2.30$\pm$1.57 & +0.33$\pm$2.09 & +0.47$\pm$0.33 \\
H06 & +0.80$\pm$1.62 & +5.69$\pm$1.31 & +0.62$\pm$0.40 & +2.10$\pm$1.00 & +1.24$\pm$0.56 & +0.03$\pm$0.38 \\
H08 & +2.04$\pm$1.65 & +1.99$\pm$1.28 & +0.58$\pm$0.85 & +4.31$\pm$1.98 & +2.21$\pm$3.24 & $-$0.06$\pm$0.61 \\
\bottomrule
\addlinespace[3pt]
\multicolumn{7}{@{}p{\dimexpr\textwidth-2\tabcolsep\relax}@{}}{Kartlar Tablo~\ref{tab:kitti-test}'dedir. Val $\Delta$: DontCare kurallı mAP50-95. İç$-$yeniden: eğitim-içi en iyi val (Ultralytics) $-$ aynı ağırlığın COCO-tarzı DontCare'siz yeniden değerlendirmesi. B satırı (taban model) mutlak değerdir; AP$_S$/AP$_M$/AP$_L$ COCO alan sınıflarıdır.} \\
\end{tabularx}
\end{table}


Tablo~\ref{tab:kitti-test-ek} kart başına ikincil test ölçütlerini, kart başına doğrulama $\Delta$'sını ve eğitim-içi ile yeniden değerlendirme farkını verir; B satırı mutlak değerdir. YOLO11s'te mAP50 ($+0.80$ ile $+2.14$), AP$_S$ ($+1.28$ ile $+5.69$), AP$_M$ ($+0.58$ ile $+1.60$) ve AP$_L$'nin ($+1.28$ ile $+4.31$) dördünde de yedi kartın yedisi pozitiftir; doğrulama $\Delta$'sı da yedi kartta pozitiftir. Bu, tek bir kartın özgül etkisinden çok, tüm kartları aynı taban modele karşı kaydıran ortak bir farkla uyumludur (Ek~\ref{sec:ek-kontrol-noktasi}). YOLO26n'de işaretler karışıktır: pozitif kart sayısı mAP50'de 4/7, AP$_S$'de 3/7, AP$_M$'de 4/7, AP$_L$'de 4/7'dir. Son sütun, eğitim-içi en iyi doğrulama değeri ile aynı \texttt{best.pt} ağırlığının COCO-tarzı \texttt{DontCare}'siz yeniden değerlendirmesi arasındaki farktır: 24 koşum ortalaması YOLO26n'de $+0.37$, YOLO11s'te $+0.40$, 48 koşum aralığı $-1.03$ ile $+1.80$'dir (37/48 koşumda pozitif). \texttt{best.pt} Ultralytics uygunluğuyla seçildiğinden, birincil uç noktanın değerlendiricisinden farklı bir ölçüte dayanır. AP$_S$ COCO'nun küçük nesne sınıfıdır ve YOLO26n'de s.d.'ler bu sütunda 1.21 ile 3.08 arasındadır; tek tek kart farkları bu gürültüyle birlikte okunmalıdır.

% uretici: evidence/scripts/build_ek_tables.py <- evidence/evidence.json (ek4.test.*.sinif.*, ek3.b2.test_nesne.*; elle duzenleme bir sonraki uretimde kaybolur)
\begin{table}[!htbp]
\centering\small\setlength{\tabcolsep}{2pt}
\caption{KITTI test: sınıf başı AP50-95 (DontCare kurallı, \texttt{best.pt}) kart $-$ taban model $\Delta$ ortalama $\pm$ s.d. [puan], $n=3$ tohum-eşli. Betimsel, renk yok.}\label{tab:kitti-sinif}
\begin{tabularx}{\textwidth}{>{\raggedright\arraybackslash}X c c c c c c c c}
\toprule
Kart & car & van & truck & ped. & sit. & cyc. & tram & misc \\
\midrule
\multicolumn{9}{@{}l}{\textbf{YOLO26n}} \\
B & 58.1 & 12.4 & 13.7 & 24.5 & 1.1 & 7.2 & 33.3 & 0.6 \\
H01 & $-$0.6$\pm$2.0 & +0.2$\pm$1.9 & $-$1.6$\pm$4.0 & +1.8$\pm$2.0 & +0.5$\pm$0.5 & +2.5$\pm$2.6 & $-$1.5$\pm$11.0 & $-$0.1$\pm$0.4 \\
H02 & +0.3$\pm$1.2 & $-$1.6$\pm$1.9 & $-$3.3$\pm$3.0 & +0.4$\pm$2.4 & $-$0.3$\pm$0.8 & +2.4$\pm$0.8 & $-$5.5$\pm$12.3 & +0.1$\pm$0.3 \\
H03 & +0.5$\pm$1.3 & +1.8$\pm$0.4 & +1.4$\pm$2.9 & +1.8$\pm$0.6 & +0.6$\pm$0.9 & +6.1$\pm$1.8 & $-$4.3$\pm$5.2 & +0.3$\pm$0.3 \\
H04 & $-$1.3$\pm$1.3 & $-$0.5$\pm$2.5 & $-$0.3$\pm$1.7 & $-$2.4$\pm$2.0 & +0.8$\pm$1.1 & +0.7$\pm$1.3 & $-$7.2$\pm$10.7 & +0.2$\pm$0.5 \\
H05 & +1.3$\pm$0.4 & +0.8$\pm$1.8 & $-$3.2$\pm$4.3 & +0.9$\pm$0.5 & $-$0.5$\pm$0.1 & +1.6$\pm$2.8 & +3.1$\pm$5.9 & +0.8$\pm$0.7 \\
H06 & $-$0.3$\pm$1.1 & $-$0.6$\pm$1.3 & $-$0.4$\pm$5.6 & $-$0.1$\pm$1.2 & $-$0.2$\pm$0.8 & +2.1$\pm$3.6 & $-$2.6$\pm$8.6 & +0.1$\pm$0.3 \\
H08 & +0.7$\pm$1.6 & $-$0.8$\pm$2.0 & +1.0$\pm$4.4 & +2.2$\pm$1.4 & +0.6$\pm$0.7 & +5.2$\pm$1.1 & $-$3.0$\pm$5.9 & $-$0.1$\pm$0.4 \\
\midrule
\multicolumn{9}{@{}l}{\textbf{YOLO11s}} \\
B & 57.7 & 12.8 & 13.5 & 27.0 & 3.4 & 15.7 & 22.9 & 1.4 \\
H01 & +2.2$\pm$0.9 & +2.4$\pm$0.9 & +0.1$\pm$3.0 & +1.8$\pm$2.4 & $-$0.1$\pm$0.7 & $-$2.6$\pm$3.4 & +8.2$\pm$3.0 & $-$0.2$\pm$0.5 \\
H02 & +1.0$\pm$1.0 & $-$0.4$\pm$3.0 & +1.6$\pm$3.2 & +1.1$\pm$0.7 & $-$0.6$\pm$0.7 & $-$1.9$\pm$4.9 & +7.8$\pm$5.6 & +0.5$\pm$0.6 \\
H03 & +1.7$\pm$1.6 & +2.8$\pm$0.5 & +3.9$\pm$3.4 & +1.9$\pm$4.2 & +0.7$\pm$5.1 & $-$3.0$\pm$3.2 & +2.9$\pm$4.1 & +0.3$\pm$1.6 \\
H04 & $-$0.6$\pm$2.6 & +3.1$\pm$5.2 & +1.1$\pm$4.5 & $-$0.3$\pm$1.0 & +2.5$\pm$4.8 & $-$1.1$\pm$4.6 & +3.5$\pm$11.8 & +0.3$\pm$1.1 \\
H05 & +4.3$\pm$1.8 & +2.2$\pm$2.5 & +4.1$\pm$2.5 & +1.7$\pm$0.5 & $-$0.8$\pm$5.4 & $-$1.8$\pm$4.3 & +8.3$\pm$6.8 & +0.2$\pm$1.1 \\
H06 & +0.0$\pm$2.1 & +1.5$\pm$3.9 & +3.3$\pm$5.1 & +2.6$\pm$2.9 & +0.6$\pm$1.2 & $-$2.1$\pm$4.9 & $-$1.1$\pm$10.9 & $-$0.1$\pm$1.0 \\
H08 & +2.1$\pm$4.0 & +2.6$\pm$1.9 & +1.4$\pm$4.3 & +4.0$\pm$2.4 & +1.9$\pm$8.5 & $-$2.2$\pm$6.3 & +2.7$\pm$21.6 & +0.9$\pm$1.3 \\
\bottomrule
\addlinespace[3pt]
\multicolumn{9}{@{}p{\dimexpr\textwidth-2\tabcolsep\relax}@{}}{Kartlar Tablo~\ref{tab:kitti-test}'dedir; sit. = person\_sitting. B satırı mutlak AP. Test nesne sayısı: car 4457, van 464, truck 119, ped. 574, sit. 56, cyc. 232, tram 126, misc 184.} \\
\end{tabularx}
\end{table}


Tablo~\ref{tab:kitti-sinif} sınıf başı AP50-95 $\Delta$'sını verir. Sınıfların test nesne sayıları car'dan çok azdır (sayılar yukarıda ve tablo dipnotunda); bu yüzden az örnekli sınıflarda s.d.'ler farkların kendisinden büyük kalabilir (örnek: YOLO11s H01'de truck $+0.1\pm3.0$) ve kazancın sınıflara dağılımı bu veriyle ayrıştırılamaz.

\subsection{H05'in KITTI'deki eğitim eğrisi}\label{sec:ek-h05-kitti}
Aynı müdahale KITTI'de VOC'taki kaybı vermedi (Bölüm~\ref{sec:genelleme-sonuc}). Ölçülen: H05, iki modelde de taban modelden daha geç tepe yapar (en iyi epoch YOLO11s'te 47/54/17, taban modelde 9/32/1; YOLO26n'de 46/39/33, taban modelde 16/11/28; 6/6 tohumda daha geç, ancak YOLO11s tohum 2'de taban modelin en iyi epoch'u 1 olduğundan bu halka zayıftır). Geç tepe en tutarlı biçimde H05'te görülür: diğer altı kartta taban modelden daha geç tepe, iki modelin altı tohumundan en çok 4'ünde görülür (H02, H08). Doğrulama eğrisinin 11.--90. epoch ortalaması farkı YOLO26n'de kartlar arasında en büyüktür (+1.29 $\pm$ 0.20, üç tohumda pozitif); YOLO11s'te ise H05'e özgü değildir (+0.53 $\pm$ 0.37; H03 +1.01, H08 +1.00; Tablo~\ref{tab:kitti-ikincil}). Mekanizma her koşumda devredeydi. Kalan hipotez ve sınırı için bkz. Bölüm~\ref{sec:genelleme-sonuc} (S2'nin yanıtı).

\subsection{Kontrol noktası seçimi: \texttt{best.pt} analizi}\label{sec:ek-kontrol-noktasi}
Bu alt bölüm, ana metinde özetlenen \texttt{best.pt} seçimine duyarlılığın ayrıntısıdır (Bölüm~\ref{sec:genelleme-sonuc}); analiz keşifseldir. Belirti, YOLO11s'te yedi kartın yedisinde pozitif çıkan \texttt{best.pt} $\Delta$'sıdır (VOC'ta aşağı yönlü olan H02 ve H05 dahil; Şekil~\ref{fig:best-last}b). Dört aday açıklama sınandı (sonuç: aşağıda): erken tepe, 10 sahne grubundan oluşan doğrulama kümesinin gürültüsü, taban model tohum 2'nin aykırılığı ve \texttt{best.pt}'nin \texttt{DontCare}'siz doğrulama uygunluğuyla seçilmesi (birincil uç nokta \texttt{DontCare} kurallıdır).

{\pgfplotsset{every axis post/.append style={height=4.2cm}}% yalnız görsel yükseklik
% F7 -- KITTI test: kart basina best.pt ve last.pt Δ mAP50-95 (3 seed ort. ± 1 s.d.), model basina bir panel (alt alta).
% Veri: figures/data/f7a_yolo26n.dat, f7b_yolo11s.dat (uretici: figures/scripts/f7_best_last.py <- evidence/evidence.json ek.test.*). Figur ici metin Turkce.
% Renk disi kanal: best.pt = dolu daire (Holm-anlamli iyilesme = dolu ucgen, mavi), last.pt = bos kare.
\begin{figure}[!htbp]
\centering
\pgfplotstableset{col sep=semicolon}
\pgfplotsset{
  fseven/.style={scale only axis, width=0.52\textwidth, xmin=-4.6, xmax=5.0, ymin=0.4, ymax=7.6, ytick=data, font=\small,
    tick label style={font=\small}, label style={font=\small}, title style={font=\small, yshift=-1pt},
    xmajorgrids=true, grid style={black!10}, axis line style={black!70}, clip=false, every axis plot/.append style={forget plot}},
  fseven best/.style={only marks, mark=*, mark size=2.4pt, draw=black, fill=black,
    /pgfplots/error bars/x dir=both, /pgfplots/error bars/x explicit, /pgfplots/error bars/error bar style={line width=1.4pt, black}, /pgfplots/error bars/error mark options={rotate=90, mark size=2.6pt, black}},
  fseven good/.style={only marks, mark=triangle*, mark size=3.2pt, draw=iyiRenk, fill=iyiRenk,
    /pgfplots/error bars/x dir=both, /pgfplots/error bars/x explicit, /pgfplots/error bars/error bar style={line width=1.4pt, iyiRenk}, /pgfplots/error bars/error mark options={rotate=90, mark size=2.6pt, iyiRenk}},
  fseven last/.style={only marks, mark=square, mark size=2.3pt, draw=black!60, fill=white, line width=0.9pt,
    /pgfplots/error bars/x dir=both, /pgfplots/error bars/x explicit, /pgfplots/error bars/error bar style={line width=1.0pt, black!60, dashed}, /pgfplots/error bars/error mark options={rotate=90, mark size=2.6pt, black!60}},
}
\begin{tikzpicture}
% ---------------- panel (a): YOLO26n
\begin{axis}[name=A, fseven, height=6.4cm, yticklabels from table={\figroot data/f7a_yolo26n.dat}{label},
  title={(a) YOLO26n, KITTI test (13 sahne grubu)}, xticklabels={}]
\draw[line width=1.1pt, black] (axis cs:0,0.4) -- (axis cs:0,7.6);
\addplot[draw=none, mark=none] table[x expr=0, y=ypos]{\figroot data/f7a_yolo26n.dat};% ytick=data icin sirali tam y listesi
\addplot[fseven best, y filter/.expression={\thisrow{isgood}==0 ? y : nan}] table[x=best_m, y expr=\thisrow{ypos}+0.17, x error=best_sd]{\figroot data/f7a_yolo26n.dat};
\addplot[fseven good, y filter/.expression={\thisrow{isgood}==1 ? y : nan}] table[x=best_m, y expr=\thisrow{ypos}+0.17, x error=best_sd]{\figroot data/f7a_yolo26n.dat};
\addplot[fseven last] table[x=last_m, y expr=\thisrow{ypos}-0.17, x error=last_sd]{\figroot data/f7a_yolo26n.dat};
\end{axis}
% ---------------- panel (b): YOLO11s
\begin{axis}[name=B, fseven, at={(A.south west)}, anchor=north west, yshift=-1.4cm, height=6.4cm, yticklabels from table={\figroot data/f7b_yolo11s.dat}{label},
  title={(b) YOLO11s, KITTI test (13 sahne grubu)}, xlabel={$\Delta$ test mAP50-95, aynı tohumlu taban modele göre [puan]},
  legend style={at={(0.5,-0.24)}, anchor=north, legend columns=2, font=\footnotesize, draw=none, /tikz/every even column/.append style={column sep=8pt}}]
\draw[line width=1.1pt, black] (axis cs:0,0.4) -- (axis cs:0,7.6);
\addplot[draw=none, mark=none] table[x expr=0, y=ypos]{\figroot data/f7b_yolo11s.dat};% ytick=data icin sirali tam y listesi
\addplot[fseven best, y filter/.expression={\thisrow{isgood}==0 ? y : nan}] table[x=best_m, y expr=\thisrow{ypos}+0.17, x error=best_sd]{\figroot data/f7b_yolo11s.dat};
\addplot[fseven good, y filter/.expression={\thisrow{isgood}==1 ? y : nan}] table[x=best_m, y expr=\thisrow{ypos}+0.17, x error=best_sd]{\figroot data/f7b_yolo11s.dat};
\addplot[fseven last] table[x=last_m, y expr=\thisrow{ypos}-0.17, x error=last_sd]{\figroot data/f7b_yolo11s.dat};
\addlegendimage{only marks, mark=*, mark size=2.4pt, black}\addlegendentry{\texttt{best.pt}, ortalama $\pm$1 s.d. (3 tohum)}
\addlegendimage{only marks, mark=square, mark size=2.3pt, black!60, fill=white, line width=0.9pt}\addlegendentry{\texttt{last.pt}, ortalama $\pm$1 s.d. (3 tohum)}
\addlegendimage{only marks, mark=triangle*, mark size=3.2pt, iyiRenk}\addlegendentry{iyi, Holm düzeltmeli $p<0.10$, $\Delta\ge+0.40$}
\end{axis}
\end{tikzpicture}
\caption{KITTI test: kart başına taban modele göre $\Delta$ mAP50-95 [puan], \texttt{best.pt} ve \texttt{last.pt} için üç tohum ortalaması $\pm$1 s.d.; (a) YOLO26n, (b) YOLO11s. Mavi üçgen Holm-5 kabulüdür.}\label{fig:best-last}
\end{figure}
}

\paragraph{Kart başına \texttt{best.pt} ve \texttt{last.pt}.}
Şekil~\ref{fig:best-last} kart başına \texttt{best.pt} ve \texttt{last.pt} $\Delta$'sını yan yana çizer. YOLO11s'te (b) \texttt{best.pt} $\Delta$'sı, VOC'ta aşağı yönlü olan H02 ve H05 dahil, yedi kartın yedisinde pozitiftir; \texttt{last.pt} $\Delta$'ları ise sıfır çevresinde dağılır. YOLO26n'de (a) örüntü aynadır: \texttt{best.pt}'de üç kart negatiftir, \texttt{last.pt}'de yedi kartın yedisi taban modelin üstündedir. Sağlamlık etiketi 14 kartın 7'sinde ``\texttt{best.pt} seçimine duyarlı'' çıktı (Tablo~\ref{tab:kitti-test}); YOLO11s'te H02 ve H05'in ``seçimden bağımsız'' etiketi, \texttt{last.pt} $\Delta$'ları $\approx$0 olduğundan işaret kuralının sınır durumudur. Kartların ortak kayması (yedi kartın ortalama test mAP50-95'i eksi taban model ortalaması) YOLO11s'te \texttt{best.pt}'de +1.38, \texttt{last.pt}'de $-$0.27 puandır (YOLO26n'de +0.01 ve +0.65), yani taban modelin kontrol noktasına bağlıdır. Yedi kart tek bir taban modele karşı ölçüldüğünden 7/7 yedi bağımsız kanıt değildir; kümeli GA'nın YOLO11s'te beş kartta sıfırı dışlaması da kart etkisinden çok bu ortak kaymayla uyumludur. Hipotez kalan açıklama: 10 sahne gruplu küçük doğrulama kümesinde \texttt{best.pt} seçimi (YOLO11s'te sık sık epoch-1 ağırlığı) bu kümeye özgü, testte kısmen taşınmayan bir tepeyi yakalar (seçilenin laneti ve sahne grubu gürültüsü). Diğer açıklamalar ve eğitim hattı tanısı Ek~\ref{sec:ek-kontrol-noktasi}'tedir; açıklamayı ayıracak deney (ara epoch'larda sabit epoch değerlendirmesi ya da daha çok tohum) koşulmadı.

\paragraph{Epoch-1 ağırlıkları.}
YOLO11s'in 24 koşumunun 11'inde en iyi epoch 1'dir, 16'sında 20'nin altındadır ve medyan 10.5'tir (Tablo~\ref{tab:kitti-epoch}); YOLO26n'de en iyi epoch 1 olan koşum yoktur (0/24), 10'u 20'nin altındadır ve medyan 24'tür. YOLO11s tohum 2'de sekiz kolun altısında (taban model, H01, H02, H03, H04, H06) \texttt{best.pt} epoch-1 ağırlığıdır. Kart--taban model kıyaslarından beşi iki epoch-1 ağırlığı arasındadır ve test $\Delta$ ortalaması +1.235'tir (aralık +0.634 ile +1.593); beş kıyas aynı taban model ağırlığına dayandığından bağımsız değildir. Taban model tohum 2'nin test mAP50-95'i \texttt{best.pt}'de 18.22, \texttt{last.pt}'de 19.83'tür.

\paragraph{Eğitim hattı.}
1000 eğitim görüntüsünde mAP50-95 şöyledir: YOLO26n taban model tohum 0'da \texttt{best.pt} 56.06, \texttt{last.pt} 78.42 (aynı \texttt{best.pt}'nin doğrulama değeri 17.89); YOLO11s taban model tohum 0--2'de \texttt{last.pt} 85.13, 85.00 ve 85.25 (aynı ağırlıkların doğrulama değeri 14.8--15.9). Epoch-1 \texttt{best.pt} eğitim sahnelerini henüz öğrenmemiştir: YOLO11s tohum 2'de 36.74 (tohum 0 ve 1'de, en iyi epoch 9 ve 32, 49.84 ve 69.93). Düşük doğrulama ve test düzeyi hat hatasıyla değil, genelleme açığıyla uyumludur. Olası nedenler 36 eğitim sahne grubuna aşırı uyum ve VOC'taki davranışı korumak için sabitlenen optimizatör ile öğrenme hızının (MuSGD, 0.01) bu veriye uymamasıdır. Üçüncü bir olasılık, tam ince ayarın önceden eğitilmiş özellikleri sahne-ayrık doğrulama ve testte bozmasıdır; bu etki büyük dağılım kayması altında gösterilmiştir \cite{kumar_2022_finetuning}, ancak KITTI bölmelerindeki kaymanın büyüklüğü ölçülmedi. Üç olasılığın hiçbiri sınanmadı.

\paragraph{Seçim ölçütü uyumsuzluğu.}
Uyumsuzluk hipotezi, geç epoch'ların \texttt{DontCare} bölgelerindeki tespitlerle cezalandırıldığını varsayar ve \texttt{DontCare} kuralıyla best$-$last farkının daralmasını öngörür. Aynı ağırlıklar bu çalışmanın COCO-tarzı değerlendiricisiyle doğrulamada \texttt{DontCare}'siz ve \texttt{DontCare} kurallı ölçüldüğünde ise fark 48 koşumun 48'inde genişler (ortalama YOLO26n'de +3.23'ten +3.62'ye, YOLO11s'te +3.84'ten +4.31'e); hipotez desteklenmez. Doğrulamada \texttt{best.pt}'nin \texttt{last.pt}'den 48/48 koşumda iyi olması yapısaldır (\texttt{best.pt} doğrulama üzerinde seçilir) ve ayırt edici değildir. Ultralytics'in eğitim-içi (\texttt{DontCare}'siz) değerlendiricisi ayrı bir ölçümdür; \texttt{best.pt} birincil uç noktanın değerlendiricisiyle değil, onun uygunluğuyla seçilir (sayılar: Tablo~\ref{tab:kitti-epoch}). Testte taban modelin \texttt{best.pt} $-$ \texttt{last.pt} farkı YOLO26n'de +1.50, YOLO11s'te $-$0.36'dır; \texttt{DontCare} kurallı doğrulama farkı ise +3.78 ve +3.43 puandır. YOLO11s'te doğrulama tepesinin testte karşılığı görülmez, YOLO26n'de kısmen taşınır.

\paragraph{Ortak kayma.}
Kartların taban modele göre ortak kayması (yedi kartın ortalama test mAP50-95'i eksi taban model ortalaması) YOLO11s'te \texttt{best.pt}'de +1.38, \texttt{last.pt}'de $-$0.27 puandır (fark 1.65). YOLO26n'de örüntü aynadır: kayma \texttt{best.pt}'de +0.01, \texttt{last.pt}'de +0.65'tir (Şekil~\ref{fig:best-last}a). Ortak kayma, taban modelin hangi kontrol noktasıyla temsil edildiğine bağlıdır; yedi kart tek bir taban modele karşı ölçüldüğünden 7/7 yedi bağımsız kanıt değildir. Tohum 2'deki aykırılık örüntüyü tek başına açıklamaz: YOLO11s'te tohum 0 ve 1'de de yedi kartın yedisi pozitiftir (tohum başına ortak kayma +1.36 ve +0.78; tohum 2'de +1.99). İşaret kuralı YOLO11s'teki H02 ve H05'e ``seçimden bağımsız'' etiketi verir; ancak bu iki kartın \texttt{last.pt} $\Delta$'ları sıfırdır (+0.00, +0.03) ve \texttt{best.pt}'deki büyüklük (+1.14, +2.28) seçime bağlıdır. Etiket bu iki kartta işaret kuralının sınır durumudur.

\paragraph{Sonuç.}
H01 kabulü önceden yazılmış kuralla raporlanır, ama ``\texttt{best.pt} seçimine duyarlı'' etiketiyle ve kazancın çoğunun tek bir test sahne grubundan geldiği notuyla; sağlam bir iyileşme iddiası yapılmaz. Ölçülen: YOLO11s'te \texttt{best.pt} ve \texttt{last.pt} $\Delta$'ları sistematik biçimde ayrışır, YOLO26n'de ayrışma ters yöndedir. Hipotez kalan: \texttt{best.pt} seçimi, küçük bir doğrulama kümesine (10 sahne grubu) özgü bir tepeyi yakalar ve bu tepe testte kısmen taşınmaz (seçim yanlılığı, ``winner's curse'', ve sahne grubu gürültüsü). Bunu ayıran deney (ara epoch'larda sabit epoch değerlendirmesi ya da daha çok tohum) koşulmadı.

\FloatBarrier
\subsection{KITTI sağlamlık analizi (kilit sonrası, keşifsel)}\label{sec:ek-saglamlik}
Bu analizler KITTI testi açıldıktan sonra yapıldı; yeni test açılışı, yeni çıkarım ya da eğitim yoktur, mevcut \texttt{best.pt} ve \texttt{last.pt} test tahminleri yeniden okundu ve tam örneklemde kayıtlı $\Delta$'lar yeniden üretildi. Önceden yazılmış kurallarda yer almazlar ve Holm düzeltmesi uygulanmadı.

% uretici: evidence/scripts/build_ek_tables.py <- evidence/evidence.json (ek5.sag.*; elle duzenleme bir sonraki uretimde kaybolur)
\begin{table}[p]
\centering\small\setlength{\tabcolsep}{3pt}
\caption{KITTI test sağlamlık (kilit sonrası, keşifsel): kart $-$ taban model $\Delta$ [puan]; tohum ve 13 sahne grubunda hiyerarşik bootstrap ($B=2000$). Mavi+kalın: GA sıfırı dışlar.}\label{tab:kitti-saglamlik}
\begin{tabularx}{\textwidth}{>{\raggedright\arraybackslash}X c c c c c c}
\toprule
Kart & $\Delta$ & Hiyerarşik GA95 & $P(\Delta\le0)$ & Tram'sız $\Delta$ & 6 sınıf $\Delta$ & Poz. grup \\
\midrule
\multicolumn{7}{@{}l}{\textbf{YOLO26n}} \\
H01 multi-scale & +0.12 & [$-$2.26, +2.51] & 0.401 & +0.36 & +0.11 & 5/13 \\
H02 Inner-CIoU & $-$0.93 & [$-$2.82, +0.85] & 0.778 & $-$0.28 & $-$1.21 & 5/13 \\
H03 kutu copy-paste & +1.01 & [$-$0.81, +3.36] & 0.093 & +1.78 & +1.20 & 7/13 \\
H04 mosaic scale & $-$1.24 & [$-$3.05, +1.65] & 0.758 & $-$0.38 & $-$1.81 & 9/13 \\
H05 Varifocal & +0.61 & [$-$1.53, +2.32] & 0.278 & +0.25 & +0.77 & 7/13 \\
H06 kutu kazancı & $-$0.25 & [$-$2.06, +1.54] & 0.600 & +0.09 & $-$0.31 & 5/13 \\
H08 yüksek doz & +0.71 & [$-$0.93, +2.73] & 0.147 & +1.24 & +0.87 & 7/13 \\
\midrule
\multicolumn{7}{@{}l}{\textbf{YOLO11s}} \\
H01 multi-scale & +1.49 & [$-$0.67, +2.61] & 0.077 & +0.53 & +2.03 & 9/13 \\
H02 Inner-CIoU & +1.14 & [$-$0.76, +2.37] & 0.120 & +0.19 & +1.53 & 10/13 \\
H03 kutu copy-paste & +1.41 & [$-$0.36, +3.50] & 0.050 & +1.20 & +1.70 & 11/13 \\
H04 mosaic scale & +1.07 & [$-$1.00, +3.01] & 0.116 & +0.72 & +0.96 & 11/13 \\
H05 Varifocal & \iyi{+2.28} & \iyi{[+0.20, +4.10]} & \iyi{0.018} & +1.42 & +3.15 & 9/13 \\
H06 kutu kazancı & +0.57 & [$-$0.81, +2.68] & 0.150 & +0.82 & +0.68 & 10/13 \\
H08 yüksek doz & +1.68 & [$-$0.45, +4.62] & 0.084 & +1.54 & +1.78 & 11/13 \\
\bottomrule
\addlinespace[3pt]
\multicolumn{7}{@{}p{\dimexpr\textwidth-2\tabcolsep\relax}@{}}{$P(\Delta\le0)$: bootstrap'ta $\Delta\le0$ oranı. Tram'sız: 7 sınıf; 6 sınıf: person\_sitting ve misc hariç; bu iki sütunda renk yoktur (GA gösterilmez). Poz. grup: tohum-ortalama $\Delta>0$ olan sahne grubu. \texttt{best.pt}. H02, H05: yönlü kartlar (beklenen $\Delta<0$).} \\
\end{tabularx}
\end{table}


Tablo~\ref{tab:kitti-saglamlik}'teki hiyerarşik bootstrap her çoğaltmada önce üç tohumu, sonra 13 test sahne grubunu yerine koyarak çeker; grup çekimi tüm kartlar ve tohumlar için ortaktır (eşli). Tohum çekimi $n=3$ üzerinden yapıldığından aralıklar hâlâ iyimser olabilir. Aralık 14 kartın 13'ünde sıfırı kapsar (YOLO11s H01 sayıları: Bölüm~\ref{sec:genelleme-sonuc}). Sekiz sınıflı $\Delta$'da sıfırı dışlayan tek kart (mavi) YOLO11s H05'tir ([+0.20, +4.10]). H02 ve H05 yönlü replikasyon kartlarıdır (beklenen $\Delta<0$) ve H05'in \texttt{last.pt} $\Delta$'sı $\approx$0'dır; bu yüzden mavi hücre bir iyileşme iddiası değildir. Tram'sız ve 6 sınıf sütunlarında aralık gösterilmediğinden renk verilmez; bu alt kümelerin aralığı tram'sız $\Delta$'da YOLO11s H05 için [+0.11, +3.32], H08 için [+0.07, +3.34]'tür. Pozitif grup sütunu görüntü ağırlıksızdır ve küçük gruplarda (5--6 görüntü) tek-grup $\Delta$'ları kararsızdır.

YOLO11s H01 kabulünün diğer duyarlılıkları şunlardır. Kümeli bootstrap'ın tek yönlü $p$'si 0.0535'tir (önceden yazılmış değil, duyarlılık). Kabul, $t$ testinin normallik varsayımına dayanır: üç tohumla işaret-çevirme testinin ulaşabileceği en küçük tek yönlü $p$ 0.125'tir. Kural belgesi modeller arası havuzlamayı öngörmez; duyarlılık olarak iki modelin 10 iyileşme sınaması birlikte Holm-10 ile düzeltildiğinde $p=0.0848$, 14 testin tümü tek ailede toplandığında $p=0.119$ çıkar. Bu değerler önceden yazılmış kararı değiştirmez. drive\_0056 dışarıda bırakıldığında H01'in $\Delta$'sı +1.49'dan +0.45'e iner ($\approx$\%70 düşüş); mAP gruplar üzerinde toplamsal olmadığından bu bir dışarıda bırakma farkıdır. Kazanç standart üç sınıfta (car, pedestrian, cyclist) görülmez, az örnekli sınıflardan gelir (Tablo~\ref{tab:kitti-ikincil}).

Kümeli bootstrap GA, sabit ağırlıklarda yalnız test sahne gruplarının örnekleme belirsizliğini ölçer ve tohum varyansını içermez. Kabul alamayan iyileşme kartlarında üst sınırı en çok +2.92 (YOLO26n) ve +2.76'dır (YOLO11s). YOLO11s'te yedi kartın beşinde (H03, H04, H05, H06, H08) sıfırı dışlar, kabul alan H01'de ise keser; YOLO26n'de hiçbir kartta sıfırı dışlamaz. Bu örüntü kart-özgü bir etkiden çok, tüm kartları aynı taban modele karşı kaydıran ortak bir farka işaret eder (Ek~\ref{sec:ek-kontrol-noktasi}) ve kart etkisinin kanıtı sayılmaz.

Doğrulama (987 görüntü, 10 sahne grubu; ikincil uç nokta) testle aynı kartı seçmedi. YOLO26n'de doğrulamada da kabul yoktur. YOLO11s'te doğrulamada kabul ölçütünü yalnız H04 sağladı (+1.82 $\pm$ 0.43, $p_{\mathrm{Holm}}=0.046$; karar değildir); H01 doğrulamada +1.73 $\pm$ 1.78 ile kabul almadı ($p_{\mathrm{Holm}}=0.350$). Testte tablo tersine döner: H01 ölçütü sağlar, H04 ($+1.07\pm0.72$, $p_{\mathrm{Holm}}=0.252$) sağlamaz. İki kümenin farklı kart seçmesi, YOLO11s'teki tek kabulün hangi kartta çıktığının sahne grubu örneklemesine duyarlı olmasıyla uyumludur.

\subsection{Eğitimsiz çıkarım (H07) ayrıntıları}\label{sec:ek-kitti-h07}
H07 taban model \texttt{best.pt} ağırlıklarıyla ve doğrulamada yapıldı; teste model başına doğrulamada en yüksek ortalamalı tek satır girdi. YOLO11s'te bu ağırlıkların tohum 2'sinin \texttt{best.pt}'si epoch-1 ağırlığıdır (Ek~\ref{sec:ek-kontrol-noktasi}). YOLO26n'de seçilen U832 satırı doğrulamada +1.50, testte +1.54 $\pm$ 0.58 puan verdi; gecikme batch 1'de 7.74 ms'den 7.80 ms'ye çıktı (oran 1.008), yani önceden yazılmış $\times1.10$ bütçesinin içindedir. Piksel sayısı 1.69 kat artarken gecikmenin yalnız \%0.8 artması, ölçümün görüntü başına Python ve ön-işleme ek yüküyle domine edildiğini düşündürür; VOC'ta aynı satırın gecikme oranı 1.14 idi. YOLO11s'te seçilen W2 satırı (T2 dallarının WBF füzyonu) doğrulamada +1.61, testte +2.02 $\pm$ 0.95 verdi; gecikme 6.82 ms'den 26.70 ms'ye ($\times3.91$) çıktı ve önceden yazılmış $\times1.10$ bütçesinin dışındadır. İki modelin gecikme motoru farklıdır (YOLO11s: VOC motoru; YOLO26n: görüntü başına \texttt{predict}); R0'da nano 7.74 ms, small 6.82 ms çıkar, bu yüzden mutlak ms modeller arasında kıyaslanmaz, yalnız model-içi oran okunur. Bu $\Delta$'lar aynı ağırlıkla ölçüldüğünden kartlar ile taban model arasındaki \texttt{best.pt} seçim farkını içermez; ancak tek satırlı, üç tohumlu ve keşifseldir, GA ve Holm verilmedi. VOC doğrulamasında 832 px girişi mAP50-95'i $-$5.19 düşürmüştü (Bölüm~\ref{sec:bulgu-gecikme}); KITTI'de aynı satırın işareti pozitiftir. Bu zıtlık için sınanmamış hipotez, yukarıdaki çözünürlük farkıdır (Ek~\ref{sec:ek-kitti-kurulum}).

\paragraph{Gecikme ölçüm hatası ve düzeltmesi.}
KITTI H07 betiğinin ilk koşumu, seçilen satırın gecikme ölçümünde CUDA bellek yetersizliğiyle düştü; hiçbir gecikme değeri yazılmadı. Görüntü yolları \texttt{predict}'e tek kaynak listesi olarak verilmişti ve Ultralytics listeyi tek batch'te işler; kural belgesindeki batch 1 sağlanmıyordu. Hata, hiçbir gecikme değeri üretilmeden ve KITTI testi açılmadan düzeltildi: YOLO11s için VOC H07 ölçüm motoru birebir, YOLO26n için görüntü başına \texttt{predict} ve GPU eşzamanlaması kullanıldı. 50 ısınma, 300 ölçüm ve FP16 değişmedi; doğruluk satırlarının tanımı ve seçim kuralı da değişmedi. Bildirilen KITTI gecikmeleri düzeltilmiş ölçümdendir; VOC gecikmeleri etkilenmedi.

\FloatBarrier
